\documentclass[journal]{IEEEtran}

\usepackage{cite}
\usepackage{amsmath}
\usepackage{booktabs}
\usepackage{url}

\title{Beyond Binary Compliance: An Explainable Multidimensional Device Trust Framework for Enterprise Access Decisions}

\author{Taqui Mohammad}

\begin{document}
\maketitle

\begin{abstract}
Enterprise access-control systems increasingly operate in environments where device
state, identity assurance, endpoint security, telemetry freshness, and threat signals
can change independently. Binary compliance provides a useful administrative signal
but may be insufficient as a standalone representation of device trust. This work
investigates an explainable multidimensional device-trust model that combines
normalized endpoint, identity, security, and risk signals with explicit safety gates.
We compare the model with a binary compliance baseline using reproducible synthetic
endpoint scenarios, including benign, degraded, stale, identity-risk, compromised,
and adversarially compliant cases. The study evaluates false-allow rates, false-deny
and step-up behavior, sensitivity to weights and thresholds, missing telemetry,
adversarial behavior, explainability, and computational scalability. Synthetic
experiments are used for controlled failure analysis and are explicitly separated from
planned external-telemetry validation. This manuscript is a work in progress; the
current synthetic results do not establish production effectiveness.
\end{abstract}

\begin{IEEEkeywords}
zero trust, device trust, endpoint security, access control, risk-adaptive access
control, endpoint management, explainable security
\end{IEEEkeywords}

\section{Introduction}
Zero Trust Architecture rejects implicit trust based on network location and calls
for explicit authentication and authorization of subjects and devices
\cite{nist800207}. Prior work has explored trust computation, risk-adaptive access
control, and continuous verification across several domains
\cite{syed2022zta,bhattarai2024ami,albomi2024service}. Recent systematic reviews
continue to identify challenges in standardization, scalability, validation, and
explainability of trust-scoring methods \cite{ruambo2026tsa}.

This paper focuses on a narrower enterprise-endpoint question: whether an explainable
multi-signal trust model can improve access decisions relative to binary compliance
under controlled and reproducible endpoint scenarios, while making uncertainty and
non-compensatory safety constraints explicit.

\section{Related Work}
\subsection{Zero Trust and Device-Aware Policy}
\subsection{Trust Scoring and Risk-Adaptive Access Control}
\subsection{Continuous Authentication and Continuous Verification}
\subsection{Research Gaps}

\section{Research Questions}
\begin{itemize}
    \item RQ1: Does multidimensional trust reduce false allows compared with binary compliance?
    \item RQ2: How sensitive are decisions to weights, thresholds, missing signals, and non-compensatory safety constraints?
    \item RQ3: Can the model remain explainable and computationally practical at enterprise scale?
\end{itemize}

\section{Model}
The initial model combines normalized positive signals and inverted risk signals using
a transparent weighted sum. Severe threat risk and critically low security coverage
act as non-compensatory safety gates. Model parameters are treated as experimental
hypotheses and are varied in sensitivity studies. A separate experimental identity
and behavior guard converts otherwise allowed decisions to a step-up state when
identity assurance is very weak or anomaly risk is high; this guard is evaluated as a
hypothesis rather than presented as a validated production threshold.

\section{Experimental Methodology}
\subsection{Synthetic Scenario Generator}
The controlled generator produces safe and unsafe endpoint-session scenarios including
healthy, policy-drift, stale, identity-risk, malware, protection-missing,
mixed-degradation, and adversarially compliant cases. Scenario prevalence and feature
ranges are explicit repository parameters. Synthetic labels support controlled failure
analysis but cannot establish external validity.

\subsection{Policy Baselines}
We compare binary compliance, the weighted multidimensional model at multiple allow
thresholds, and the experimental identity/behavior guard. We additionally evaluate a
literature-inspired attribute-quorum comparator: an ordinary ALLOW requires both the
aggregate 0.75 trust threshold and a minimum number of individual safety-oriented
factors to exceed a uniform factor threshold. This follows the general decision
structure of prior trust-score ZTA work that combines aggregate and per-attribute
threshold requirements \cite{bhattarai2024ami}; it is not a reproduction of that
paper's domain-specific implementation. All policies are evaluated on identical
generated rows so comparisons are paired.

\subsection{Adversarially Compliant Endpoints}
The adversarially compliant scenario deliberately preserves a high compliance signal
while degrading other trust dimensions. It is intended to test compensation failure
modes, not to estimate the prevalence of such states in production environments.

\subsection{Metrics and Statistical Analysis}
Primary operational metrics include the false-allow rate among sessions labeled unsafe
for ordinary access and the step-up rate among sessions labeled safe. We report 95\%
Wilson score intervals for binomial rates. For paired row-level ordinary-access
correctness, we use continuity-corrected McNemar comparisons because each policy is
applied to the same fixed evaluation instances. Dietterich identified McNemar's test as
an appropriate low-Type-I-error comparison for paired classifier outputs on a single
test set \cite{dietterich1998statistical}. Statistical significance is treated as
secondary to effect size and security/friction trade-offs. These analyses are
exploratory because model development used the same synthetic design; they are not
presented as preregistered confirmatory evidence.

\subsection{Sensitivity and Scale Tests}
Threshold, weight, missing-telemetry, non-compensation, and attribute-quorum studies
are evaluated separately so that changes in security outcomes can be distinguished
from changes in user friction and abstention behavior.

\section{Preliminary Synthetic Results}
Table~\ref{tab:synthetic-policy} reports the seeded 50,000-row synthetic policy
comparison. These values characterize behavior under the generator assumptions only.
They must not be interpreted as production false-allow rates.

\begin{table}[ht]
\centering
\caption{Synthetic policy comparison (50,000 rows, seed 20260903).}
\label{tab:synthetic-policy}
\begin{tabular}{lccc}
\toprule
Policy & False ALLOW & Safe STEP\_UP & Accuracy \\
\midrule
Binary compliance & 93.74\% & 0.00\% & 52.45\% \\
Additive 0.75 & 12.50\% & 5.40\% & 91.62\% \\
Additive 0.80 & 0.34\% & 13.58\% & 91.97\% \\
Guarded 0.75 & 1.24\% & 5.40\% & 96.34\% \\
Quorum 0.75, $q$=0.70, $k$=7 & 0.00\% & 6.64\% & 96.14\% \\
\bottomrule
\end{tabular}
\end{table}

The threshold increase from 0.75 to 0.80 substantially reduces synthetic false allows
but also increases safe step-up decisions. The guarded policy changes a different
region of the decision surface: under the present generator it reduces false allows
without increasing safe step-up relative to the 0.75 additive model. The exploratory
7-of-8 attribute-quorum baseline removes all false ALLOWs in this synthetic population
while increasing safe step-up to 6.64\%. Its apparent performance is strongly aligned
to the current generator: unsafe identity/adversarial scenarios generally contain at
least two weak factors, whereas healthy cases are strong across nearly all factors.
The quorum result is therefore evidence of a useful alternative non-compensatory
structure, not evidence that this setting is optimal or externally superior.
Accordingly, no production policy recommendation is made from these results.

\section{Threats to Validity}
Synthetic data cannot establish real-world security effectiveness. Construct validity,
distribution assumptions, scenario prevalence, threshold subjectivity, and domain
transfer are explicit limitations. The synthetic generator also informed iterative
policy development, creating design-overfitting risk. Confidence intervals quantify
sampling variability conditional on the generator rather than uncertainty about the
real enterprise population. The ordinary-access accuracy metric treats both STEP\_UP
and DENY as blocked access and therefore must be interpreted alongside separate
friction and false-allow metrics. Alternative policies such as the attribute-quorum
baseline can also become implicitly tuned to generator geometry; favorable synthetic
performance must therefore be challenged under benign overlap and frozen external
telemetry. External telemetry must be evaluated on temporally held-out data with
frozen proxy definitions and policy settings before stronger claims are considered.

\section{Discussion}

\section{Conclusion}

\bibliographystyle{IEEEtran}
\bibliography{references}

\end{document}
